\chapter{Étape 1 — Lire le scanner}
\label{chap:lecture}

\begin{chapeau}
Cette étape transforme un dossier de fichiers DICOM en un volume 3D en unités Hounsfield, et calcule la matrice qui
donne la position de chaque voxel en millimètres. Elle fait les mêmes contrôles que le module DICOM de 3D Slicer.
\end{chapeau}

\section{Trouver les séries}

Un dossier apporté par le chirurgien peut contenir plusieurs examens, plusieurs séries, des sous-dossiers, des
fichiers sans extension et des fichiers qui ne sont pas des images (rapports, \fichier{DICOMDIR}). Le logiciel :
\begin{enumerate}
  \item parcourt tout le dossier, sous-dossiers compris, et garde les fichiers dont l'en-tête est DICOM ;
  \item lit seulement l'en-tête (pas les pixels), ce qui est rapide ;
  \item écarte ce qui n'est pas une image en coupe (pas de position ou d'orientation) ;
  \item regroupe les fichiers par série (\code{SeriesInstanceUID}) et donne à chaque série un nom de dossier sûr
  (minuscules, chiffres, tirets), tiré du nom du dossier.
\end{enumerate}

\section{Ranger les coupes dans le bon ordre}

Les fichiers d'une série ne sont pas forcément nommés dans l'ordre, et le champ \code{InstanceNumber}, qui donne
le numéro de l'image, n'est pas fiable. Le logiciel range donc les coupes \textbf{par leur position dans l'espace}. Soient $\vb{r}$ et
$\vb{c}$ les directions des lignes et des colonnes de l'image (les deux triplets de \code{ImageOrientationPatient}).
La normale aux coupes est
\begin{equation}
  \vb{n} = \vb{r} \times \vb{c}.
\end{equation}
La coupe $k$, dont le premier pixel est en $\vb{P}_k$ (\code{ImagePositionPatient}), est à la distance
$d_k = \vb{P}_k \cdot \vb{n}$ le long de cette normale. Trier les coupes, c'est trier les $d_k$.

\section{Contrôler la géométrie}

Avant de lire les pixels, plusieurs contrôles évitent un volume faux :
\begin{itemize}
  \item \textbf{une seule orientation} : les coupes dont l'orientation diffère de la plus fréquente (à $10^{-3}$ près)
  sont écartées. Ce sont en général les images de repérage, comme le topogramme (\cref{sec:topogramme}) ;
  \item \textbf{pas de doublon} : deux coupes à la même position, la seconde est écartée ;
  \item \textbf{espacement régulier} : le pas $\Delta = \operatorname{médiane}(d_{k+1} - d_k)$ ; si un écart s'en
  éloigne de plus de 5~\%, il manque sans doute des coupes ;
  \item \textbf{épaisseur} : au-delà de 1~mm, un avertissement est donné, car la précision des surfaces en souffre.
\end{itemize}

Si l'espacement est irrégulier, le volume est rééchantillonné sur une grille régulière de pas $\Delta$, par
interpolation linéaire entre les deux coupes voisines :
\begin{equation}
  I(d) = (1 - w)\, I_k + w\, I_{k+1},
  \qquad
  w = \frac{d - d_k}{d_{k+1} - d_k},
  \qquad d_k \le d \le d_{k+1}.
\end{equation}

\section{Convertir en unités Hounsfield}

Chaque pixel stocké (\emph{stored value}, SV) devient une valeur en HU par la règle donnée dans l'en-tête :
\begin{equation}
  \tcbhighmath{\mathrm{HU} = \mathrm{SV} \times \texttt{RescaleSlope} + \texttt{RescaleIntercept}.}
\end{equation}
Le volume est gardé en entiers de 16 bits (de $-32\,768$ à $32\,767$), ce qui suffit et divise la mémoire par deux par
rapport à des réels. Si pydicom ne sait pas décoder une image compressée, le logiciel la relit avec SimpleITK.

\section{De l'indice de voxel aux millimètres}
\label{sec:affine}

C'est le calcul le plus important de l'étape : tout ce qui suit (STL, repères, guides) en dépend. Le voxel d'indices
$(i, j, k)$ (colonne, ligne, coupe) est au point
\begin{equation}
  \begin{pmatrix} x \\ y \\ z \\ 1 \end{pmatrix}
  = \underbrace{\begin{pmatrix}
    \Delta_c\, r_x & \Delta_r\, c_x & \delta_x & P_{0,x} \\
    \Delta_c\, r_y & \Delta_r\, c_y & \delta_y & P_{0,y} \\
    \Delta_c\, r_z & \Delta_r\, c_z & \delta_z & P_{0,z} \\
    0 & 0 & 0 & 1
  \end{pmatrix}}_{\mt{A}}
  \begin{pmatrix} i \\ j \\ k \\ 1 \end{pmatrix},
  \qquad
  \boldsymbol{\delta} = \frac{\vb{P}_{N-1} - \vb{P}_0}{N - 1},
  \label{eq:affine}
\end{equation}
où $\Delta_c$ est l'écart entre deux colonnes et $\Delta_r$ l'écart entre deux lignes (\code{PixelSpacing} donne
$[\Delta_r, \Delta_c]$, dans cet ordre), $\vb{P}_0$ la position du premier voxel et $\boldsymbol{\delta}$ le pas
réel d'une coupe à la suivante, calculé avec la première et la dernière coupe.

\begin{pourquoi}
On pourrait prendre $\boldsymbol{\delta} = \Delta\,\vb{n}$. Mais si le statif du scanner était incliné
(\emph{gantry tilt}), les coupes glissent les unes par rapport aux autres et le pas n'est pas parallèle à la normale.
Avec $\boldsymbol{\delta}$ mesuré sur les positions, la matrice reste exacte. L'inclinaison est
$\arccos\bigl(\abs{\boldsymbol{\delta}\cdot\vb{n}} / \norm{\boldsymbol{\delta}}\bigr)$ ; au-delà de 0,5\degre, le
logiciel le signale.
\end{pourquoi}

L'opération inverse, des millimètres aux indices, résout le système linéaire :
$(i, j, k)^{\mathsf T} = \mt{A}_{3\times3}^{-1}\,(\vb{p} - \vb{P}_0)$. Les indices obtenus peuvent être
fractionnaires : on interpole alors entre les voxels.

\begin{exemple}
Série genou--cheville du lévrier : 748 coupes de $512 \times 512$ pixels, voxel de
$0{,}625 \times 0{,}56 \times 0{,}56$~mm, soit un volume de $467 \times 287 \times 287$~mm et
196 millions de voxels (392~Mo en entiers de 16 bits). Statif non incliné.
\end{exemple}

\begin{danslecode}
\fichier{code_source/etape1_lecture_dicom.py} :
\code{inventaire()} trouve et regroupe les séries ; \code{charger()} trie, contrôle, lit, convertit en HU et
construit la matrice $\mt{A}$ (classe \code{Volume}, avec \code{vers_lps()} et \code{vers_indices()}) ;
\code{lire_pixels()} applique \code{RescaleSlope} et \code{RescaleIntercept}. Réglages dans
\fichier{parametres.py} : \code{EPAISSEUR_MAX_MM}, \code{TOLERANCE_ESPACEMENT}, \code{TOLERANCE_ORIENTATION}.
\end{danslecode}

\begin{attention}[title={Données personnelles}]
L'en-tête DICOM contient souvent le nom de l'animal, parfois celui du propriétaire, et un identifiant. L'inventaire
les lit, et le rapport de l'étape 5 les recopie (\cref{sec:securite}). Il vaut mieux rendre les scanners anonymes avant
de les importer.
\end{attention}
